\documentclass[11pt,oneside]{book}
\usepackage[paperwidth=7in,paperheight=9.19in,margin=1in]{geometry}
\usepackage{fancyhdr}
\usepackage{amsmath}
\usepackage[pdfpagelabels=PAGELABELS]{hyperref}
\pagestyle{fancy}\fancyhf{}
\fancyfoot[R]{\rightmark\ $|$\ \thepage}
\renewcommand{\headrulewidth}{0pt}
\fancypagestyle{plain}{\fancyhf{}}
\begin{document}
\frontmatter
\chapter{Preface}
This synthetic book is a test fixture for the locator spike. It has roman front matter, a two-chapter body, footers in the form page bar chapter, and one page whose footer lies.
\newpage Second preface page with text so that roman numbering has more than one page.
\newpage Third preface page.
\mainmatter
\chapter{Alpha Topics}
\section{First Section}
The alpha principle states that every span has a stable identifier derived from its page and position.
\newpage
The second alpha page discusses hyphen\-ation and ligatures such as final and flow, and “smart quotes”.
See ``Beta Basics'' on page 5 for the next idea.
\newpage
\thispagestyle{empty}
This page has no footer at all and must be interpolated.
\newpage
\fancyfoot[R]{\rightmark\ $|$\ 99}
This page carries a lying footer number that the monotonic fit must correct.
\newpage
\fancyfoot[R]{\rightmark\ $|$\ \thepage}
\chapter{Beta Basics}
\section{Equations}
The energy relation is
\begin{equation} E = mc^2 \end{equation}
and the beta rule follows from it.
\newpage
Last body page of the fixture.
\end{document}
